\documentclass[10pt,a4paper]{amsart}
\usepackage{amsmath,amssymb,amsthm,graphicx}
\usepackage{verbatim}
\newenvironment{reference}{\comment}{\endcomment}
%\newenvironment{reference}{}{}
\newenvironment{slogan}{\comment}{\endcomment}
\newenvironment{history}{\comment}{\endcomment}

% For commutative diagrams we use Xy-pic
\usepackage[all]{xy}

% We use 2cell for 2-commutative diagrams.
\xyoption{2cell}
% Theorem environments.
%
\theoremstyle{plain}
\newtheorem{theorem}[subsection]{Theorem}
\newtheorem{proposition}[subsection]{Proposition}
\newtheorem{lemma}[subsection]{Lemma}

\theoremstyle{definition}
\newtheorem{definition}[subsection]{Definition}
\newtheorem{example}[subsection]{Example}
\newtheorem{exercise}[subsection]{Exercise}
\newtheorem{situation}[subsection]{Situation}

\theoremstyle{remark}
\newtheorem{remark}[subsection]{Remark}
\newtheorem{remarks}[subsection]{Remarks}

\numberwithin{equation}{subsection}

% Macros
%
\def\lim{\mathop{\mathrm{lim}}\nolimits}
\def\colim{\mathop{\mathrm{colim}}\nolimits}
\def\Spec{\mathop{\mathrm{Spec}}}
\def\Hom{\mathop{\mathrm{Hom}}\nolimits}
\def\Ext{\mathop{\mathrm{Ext}}\nolimits}
\def\SheafHom{\mathop{\mathcal{H}\!\mathit{om}}\nolimits}
\def\SheafExt{\mathop{\mathcal{E}\!\mathit{xt}}\nolimits}
\def\Sch{\mathit{Sch}}
\def\Mor{\mathop{\mathrm{Mor}}\nolimits}
\def\Ob{\mathop{\mathrm{Ob}}\nolimits}
\def\Sh{\mathop{\mathit{Sh}}\nolimits}
\def\NL{\mathop{N\!L}\nolimits}
\def\CH{\mathop{\mathrm{CH}}\nolimits}
\def\proetale{{pro\text{-}\acute{e}tale}}
\def\etale{{\acute{e}tale}}
\def\QCoh{\mathit{QCoh}}
\def\Ker{\mathop{\mathrm{Ker}}}
\def\Im{\mathop{\mathrm{Im}}}
\def\Coker{\mathop{\mathrm{Coker}}}
\def\Coim{\mathop{\mathrm{Coim}}}

% Boxtimes
%
\DeclareMathSymbol{\boxtimes}{\mathbin}{AMSa}{"02}

%
% Macros for moduli stacks/spaces
%
\def\QCohstack{\mathcal{QC}\!\mathit{oh}}
\def\Cohstack{\mathcal{C}\!\mathit{oh}}
\def\Spacesstack{\mathcal{S}\!\mathit{paces}}
\def\Quotfunctor{\mathrm{Quot}}
\def\Hilbfunctor{\mathrm{Hilb}}
\def\Curvesstack{\mathcal{C}\!\mathit{urves}}
\def\Polarizedstack{\mathcal{P}\!\mathit{olarized}}
\def\Complexesstack{\mathcal{C}\!\mathit{omplexes}}
% \Pic is the operator that assigns to X its picard group, usage \Pic(X)
% \Picardstack_{X/B} denotes the Picard stack of X over B
% \Picardfunctor_{X/B} denotes the Picard functor of X over B
\def\Pic{\mathop{\mathrm{Pic}}\nolimits}
\def\Picardstack{\mathcal{P}\!\mathit{ic}}
\def\Picardfunctor{\mathrm{Pic}}
\def\Deformationcategory{\mathcal{D}\!\mathit{ef}}
\usepackage{cleveref}
\crefname{originallemma}{original lemma}{original lemmas}
\Crefname{originallemma}{Original lemma}{Original lemmas}
\makeatletter
\newlabel{lemma-composition-finite-type}{{03XG}{source}{}{sourcepointer-1}{}}
\newlabel{lemma-composition-finite-type@cref}{{[originallemma][0][]03XG}{[1][0][]source}{}{}{}}
\makeatother
\crefname{originallemma}{original lemma}{original lemmas}
\Crefname{originallemma}{Original lemma}{Original lemmas}
\makeatletter
\newlabel{morphisms-lemma-etale-locally-finite-presentation}{{02GR}{source}{}{sourcepointer-2}{}}
\newlabel{morphisms-lemma-etale-locally-finite-presentation@cref}{{[originallemma][0][]02GR}{[1][0][]source}{}{}{}}
\makeatother
\crefname{originallemma}{original lemma}{original lemmas}
\Crefname{originallemma}{Original lemma}{Original lemmas}
\makeatletter
\newlabel{morphisms-lemma-finite-presentation-finite-type}{{01TW}{source}{}{sourcepointer-3}{}}
\newlabel{morphisms-lemma-finite-presentation-finite-type@cref}{{[originallemma][0][]01TW}{[1][0][]source}{}{}{}}
\makeatother
\crefname{originallemma}{original lemma}{original lemmas}
\Crefname{originallemma}{Original lemma}{Original lemmas}
\makeatletter
\newlabel{morphisms-lemma-finite-type-points-morphism}{{02J3}{source}{}{sourcepointer-4}{}}
\newlabel{morphisms-lemma-finite-type-points-morphism@cref}{{[originallemma][0][]02J3}{[1][0][]source}{}{}{}}
\makeatother
\crefname{originallemma}{original lemma}{original lemmas}
\Crefname{originallemma}{Original lemma}{Original lemmas}
\makeatletter
\newlabel{morphisms-lemma-identify-finite-type-points}{{02J2}{source}{}{sourcepointer-5}{}}
\newlabel{morphisms-lemma-identify-finite-type-points@cref}{{[originallemma][0][]02J2}{[1][0][]source}{}{}{}}
\makeatother
\crefname{originallemma}{original lemma}{original lemmas}
\Crefname{originallemma}{Original lemma}{Original lemmas}
\makeatletter
\newlabel{spaces-lemma-representable-transformations-property-implication}{{02YO}{source}{}{sourcepointer-6}{}}
\newlabel{spaces-lemma-representable-transformations-property-implication@cref}{{[originallemma][0][]02YO}{[1][0][]source}{}{}{}}
\makeatother
\title{Points of finite type on an algebraic space characterized by etale presentations}
\author{The Stacks Project Authors}
\date{stacks-project:spaces-morphisms.tex:local-snapshot-20261012; GNU FDL 1.2 or later}
\begin{document}
\maketitle
\setcounter{subsection}{1}
\noindent\textit{Editorial note.} All delivered mathematical text is the author's own, copied from the pinned Stacks Project chapter file spaces-morphisms.tex; no statement or proof step was rewritten, reordered or removed.\par\medskip
\noindent\textit{Editorial note.} The proof uses the notion of a finite type point of a scheme, which is defined in the Morphisms chapter; the shrinking step at the end is justified by the cited identification of finite type points with closed points after shrinking.\par\medskip
\noindent\textit{Editorial note.} The six references of the counted spans are recorded as non-load-bearing navigation pointers; the argument that consumes them stays inside the counted proof span.\par\medskip
\noindent\textit{Editorial note.} Preamble: the three native excerpt spans of sources/stacks-spaces-morphisms/preamble.tex (lines 8-13, 14-18 and 154-217). The chapter file itself carries no preamble text.\par\medskip
\noindent\textit{Allowed standard prerequisites (editorial).}\begin{itemize}
\item The general principle translating representable properties into scheme properties. Spaces, Lemma spaces-lemma-representable-transformations-property-implication, used at proof lines 5035-5037. Makes U -> X locally of finite type.
\item Etale morphisms are locally of finite presentation, hence locally of finite type. Morphisms, Lemmas morphisms-lemma-etale-locally-finite-presentation and morphisms-lemma-finite-presentation-finite-type, used at proof lines 5038-5040. Gives local finite typeness of the etale morphisms involved.
\item Permanence of local finite typeness under composition. Lemma lemma-composition-finite-type, used at proof lines 5040-5042. Produces the locally finite type representative in the forward direction.
\item Finite type points of schemes and the identification of finite type with closed after shrinking. Morphisms, Lemmas morphisms-lemma-identify-finite-type-points and morphisms-lemma-finite-type-points-morphism, used at proof lines 5048-5056. Provides the finite type point v and turns its image into a closed point.
\end{itemize}\medskip
\ifdefined\hypertarget\hypertarget{sourcepointer-1}{}\fi
\noindent\textit{Source reference.} Original lemma 03XG: focused-seeds/AGP/sources/stacks-snapshot-20261011/tags.txt, source lines 5147--5147. Navigation pointer into the Stacks Project; the standard fact it supplies is applied inside the counted proof span and the argument resting on it stays there.\par\medskip
\ifdefined\hypertarget\hypertarget{sourcepointer-2}{}\fi
\noindent\textit{Source reference.} Original lemma 02GR: focused-seeds/AGP/sources/stacks-snapshot-20261011/tags.txt, source lines 3238--3238. Navigation pointer into the Stacks Project; the standard fact it supplies is applied inside the counted proof span and the argument resting on it stays there.\par\medskip
\ifdefined\hypertarget\hypertarget{sourcepointer-3}{}\fi
\noindent\textit{Source reference.} Original lemma 01TW: focused-seeds/AGP/sources/stacks-snapshot-20261011/tags.txt, source lines 2389--2389. Navigation pointer into the Stacks Project; the standard fact it supplies is applied inside the counted proof span and the argument resting on it stays there.\par\medskip
\ifdefined\hypertarget\hypertarget{sourcepointer-4}{}\fi
\noindent\textit{Source reference.} Original lemma 02J3: focused-seeds/AGP/sources/stacks-snapshot-20261011/tags.txt, source lines 3322--3322. Navigation pointer into the Stacks Project; the standard fact it supplies is applied inside the counted proof span and the argument resting on it stays there.\par\medskip
\ifdefined\hypertarget\hypertarget{sourcepointer-5}{}\fi
\noindent\textit{Source reference.} Original lemma 02J2: focused-seeds/AGP/sources/stacks-snapshot-20261011/tags.txt, source lines 3321--3321. Navigation pointer into the Stacks Project; the standard fact it supplies is applied inside the counted proof span and the argument resting on it stays there.\par\medskip
\ifdefined\hypertarget\hypertarget{sourcepointer-6}{}\fi
\noindent\textit{Source reference.} Original lemma 02YO: focused-seeds/AGP/sources/stacks-snapshot-20261011/tags.txt, source lines 3883--3883. Navigation pointer into the Stacks Project; the standard fact it supplies is applied inside the counted proof span and the argument resting on it stays there.\par\medskip
\begin{lemma}
\label{lemma-point-finite-type}
Let $S$ be a scheme. Let $X$ be an algebraic space over $S$.
Let $x \in |X|$. The following are equivalent:
\begin{enumerate}
\item There exists a morphism $\Spec(k) \to X$
which is locally of finite type and represents $x$.
\item There exists a scheme $U$, a closed point $u \in U$, and an \'etale
morphism $\varphi : U \to X$ such that $\varphi(u) = x$.
\end{enumerate}
\end{lemma}

\begin{proof}
Let $u \in U$ and $U \to X$ be as in (2). Then
$\Spec(\kappa(u)) \to U$ is of finite type, and $U \to X$ is
representable and locally of finite type (by the general principle
Spaces, Lemma
\ref{spaces-lemma-representable-transformations-property-implication}
and
Morphisms, Lemmas \ref{morphisms-lemma-etale-locally-finite-presentation} and
\ref{morphisms-lemma-finite-presentation-finite-type}). Hence we see
(1) holds by
Lemma \ref{lemma-composition-finite-type}.

\medskip\noindent
Conversely, assume $\Spec(k) \to X$ is locally of finite type and
represents $x$. Let $U \to X$ be a surjective \'etale morphism where $U$
is a scheme. By assumption $U \times_X \Spec(k) \to U$ is locally
of finite type. Pick a finite type point $v$ of
$U \times_X \Spec(k)$ (there exists at least one, see
Morphisms, Lemma \ref{morphisms-lemma-identify-finite-type-points}).
By
Morphisms, Lemma \ref{morphisms-lemma-finite-type-points-morphism}
the image $u \in U$ of $v$ is a finite type point of $U$.
Hence by
Morphisms, Lemma \ref{morphisms-lemma-identify-finite-type-points}
after shrinking $U$ we may assume that $u$ is a closed point of $U$, i.e.,
(2) holds.
\end{proof}

\end{document}
